\section{The original whole-cover degree and trace normalization}
\label{sec:kn-nadel-whole-cover-degree}

\begin{lemma}[Degree of the entire varying-root quotient]
\label{lem:kn-nadel-whole-cover-degree}
\lean{EnvelopeProperty.Vanishing.NadelWholeCoverDegree.degree,EnvelopeProperty.Vanishing.NadelWholeCoverDegree.degree_pos,EnvelopeProperty.Vanishing.NadelWholeCoverDegree.basisIndex_card,EnvelopeProperty.Vanishing.NadelWholeCoverDegree.finrank_eq_degree}
\leanok
Let $R$ be a nonzero commutative ring, let $N_i>0$ be the original
root degrees, and let $B=R[z_i]/(z_i^{N_i}-c_i)$ be the whole original
quotient. Put $D=\prod_iN_i$. Its actual original monic basis has
cardinality $D$, and $\operatorname{rank}_R B=D>0$.
\end{lemma}
\begin{proof}
The actual simultaneous quotient is isomorphic over $R$ to the original
successive monic-root tower. At the first stage the basis index is a
singleton; adjoining the $i$-th root multiplies it by the entire set
$\{0,\ldots,N_i-1\}$. Induction therefore gives cardinality
$\prod_iN_i$. Transport the actual tower basis through the proved whole
algebra isomorphism. The dimension theorem for bases over a nonzero
commutative ring identifies its cardinality with the module rank.
Every factor is positive, including the empty product $1$, so $D>0$.
\uses{lem:kn-adapted-varying-finite}
\end{proof}

\begin{lemma}[Original trace after arbitrary coefficient extension]
\label{lem:kn-nadel-whole-cover-trace-degree}
\lean{EnvelopeProperty.Vanishing.NadelWholeCoverDegree.trace_algebraMap,EnvelopeProperty.Vanishing.NadelWholeCoverDegree.trace_one,EnvelopeProperty.Vanishing.NadelWholeCoverDegree.baseChange_trace_scalar,EnvelopeProperty.Vanishing.NadelWholeCoverDegree.baseChange_trace_one,EnvelopeProperty.Vanishing.NadelWholeCoverDegree.degree_isUnit}
\leanok
For the preceding whole algebra,
$\operatorname{Tr}_{B/R}(r)=Dr$ for every original scalar $r$.
For any commutative $R$-algebra $S$, the entire algebra
$C=S\otimes_RB$ satisfies
$\operatorname{Tr}_{C/S}(s)=Ds$. In particular the two traces of $1$
are the original degree $D$. If $R$ is an algebra over a
characteristic-zero field, $D$ is a unit in $R$.
\end{lemma}
\begin{proof}
In the genuine original basis, multiplication by $r$ is the diagonal
matrix with every diagonal entry $r$. Its matrix trace is $Dr$.
Base change carries that same basis to an $S$-basis of the whole
$S\otimes_RB$, with the same original index set. Multiplication by
$s$ consequently has trace $Ds$ in this basis as well. The latter
trace is relative to $S$: it introduces no rank of $S$ over $R$.
The positive integer $D$ is nonzero in the original characteristic-zero
field, hence invertible there. Its image under the actual field algebra
map is a unit in $R$.
\uses{lem:kn-nadel-whole-cover-degree}
\end{proof}
